\section{Discussion}
\label{sec:discussion}

\subsection{Each corruption traces to a physical sensor failure}
\label{sec:discussion-physical}
\advisor{point 3 -- relate the choice of corruptions to physical issues: which physical/sensor cause produces which corruption.}

The six corruptions are not arbitrary perturbations; each models a documented event-camera
failure, which is what makes the benchmark a robustness test rather than a synthetic
curiosity. Table~\ref{tab:corr-taxonomy} (Section~\ref{sec:corruptions}) lists the mapping;
we expand the physical reasoning here.

\begin{itemize}
  \item \textbf{\texttt{hot\_pixel} $\leftarrow$ defective / stuck pixels and fixed-pattern
        noise.} Fabrication defects and bias mismatch make individual pixels fire
        independently of the scene, injecting a constant DC offset at fixed locations -- the
        event-camera analogue of stuck CCD/CMOS pixels.
  \item \textbf{\texttt{event\_flood} $\leftarrow$ electromagnetic interference, LED/PWM
        flicker, and refractory-limited saturation.} Strong illumination transients or EMI
        drive bursts of spurious events in a region for a short window, locally saturating the
        readout bandwidth.
  \item \textbf{\texttt{temporal\_jitter} $\leftarrow$ clock noise and timestamp error.}
        Arbitration delay, bus contention, and finite timestamp resolution misassign events
        to neighbouring time bins, perturbing temporal order while preserving marginal counts.
  \item \textbf{\texttt{event\_rate\_shift} $\leftarrow$ contrast-threshold / biasing drift
        and domain gap.} The per-pixel contrast threshold and bias currents set the event
        rate; temperature drift or a low- vs.\ high-texture (slow- vs.\ fast-motion) domain
        gap multiplies the global rate.
  \item \textbf{\texttt{polarity\_flip} $\leftarrow$ ON/OFF misclassification.} Comparator
        offset and asymmetric ON/OFF thresholds (a known DVS non-ideality) flip the
        brightness-change sign for a fraction of events.
  \item \textbf{\texttt{spatial\_dropout} $\leftarrow$ dead-pixel clusters and occlusion.}
        Manufacturing dead regions, lens contamination, or partial occlusion zero the event
        support over contiguous areas.
\end{itemize}

The value of the closed-form theory (Section~\ref{sec:theory}) is that each \emph{physical}
axis maps to a distinct \emph{statistical} axis of the membrane moments, so a detector's
behaviour on a corruption is predictable from the sensor failure it models, not measured by
trial and error.

\subsection{What biology does about a spike flood -- and why our neuron does not}
\label{sec:discussion-bio}
\advisor{point 4 -- since SNNs are bio-plausible, check neuroscience: what happens during an event flood / burst of spikes, and how do neurons adapt?}

\texttt{event\_flood} is the corruption our membrane detector handles worst per-frame
(Table~\ref{tab:mdd-frame}): a uniform rate increase inflates every channel's mean
proportionally, mimicking a genuinely busy scene (Section~\ref{sec:theory}). It is therefore
worth asking how \emph{biological} neurons -- which the SNN abstracts -- cope with sustained
high-rate input, because the answer both explains our failure and points to a remedy.

Real neurons do not integrate a flood passively. Two mechanisms dominate.
\textbf{Spike-frequency adaptation (SFA)}: in response to the onset of a constant, strong
input a neuron's firing rate decays over tens to hundreds of milliseconds, driven by
calcium-dependent and slow voltage-dependent K$^+$ currents and Na$^+$-channel
inactivation~\cite{benda2003sfa,gabbiani2014sfa}. \textbf{Short-term synaptic depression}:
heavily used synapses transiently weaken, implementing a divisive \emph{gain
control}~\cite{abbott1997depression,rothman2009gaincontrol}. The combined effect is that the
neuron normalises sustained drive and responds chiefly to \emph{abrupt changes} rather than
to the absolute input level -- precisely the invariance that would make a flood (a sustained,
structureless rate increase) detectable as anomalous, because an adapting neuron's response to
a flood diverges from its response to genuine scene structure.

The PLIF neuron used here has \emph{none} of this: a single learned, static leak $\tau$ and a
linear leaky accumulator (Eq.~\ref{eq:plif}) with no adaptation current and no synaptic
depression. So a flood simply scales the accumulator, and the per-frame membrane statistics
move along the same direction a busy clean scene would -- exactly the near-chance result we
report. This is, to our reading, why the corruption that biology is best equipped to handle is
the one our biologically-\emph{simplified} model handles worst.

This suggests a concrete, biologically-grounded line of future work
(cf.\ Section~\ref{sec:motivation}): equip the front-end with \textbf{adaptive} neurons
(an SFA current or an adapting/AdLIF unit) or a short-term-depression input stage. Under
adaptation, a sustained flood would be attenuated relative to transient scene events, so the
membrane signature of a flood would separate from that of a busy scene -- potentially both
improving \emph{detection} of \texttt{event\_flood} and \emph{restoring} downstream mAP. We
flag this as a promising direction rather than a result, keeping this paper scoped to
detection with the standard PLIF model.
